\subsection*{CU-03: Recuperar contraseña}
\addcontentsline{toc}{subsection}{CU-03: Recuperar contraseña}
\label{cu-03}

\begin{fichacu}{CU-03}{Recuperar contraseña}
  \crudFila{Responsable}{José Eduardo Prior Hernández}
  \crudFila{Actor principal}{Jugador}
  \crudFila{Actores de sistema}{Servidor de partidas, Servicio de correo, Base de datos}
  \crudFila{Prototipos}{9a, 9b}
  \crudFila{Objetivo}{Devolver el acceso a un Jugador que olvidó su contraseña, sin intervención humana y sin que el proceso revele qué correos están registrados.}
  \crudFila{Disparador}{El Jugador selecciona ``Olvidé mi contraseña'' en la \texttt{GUI\_\allowbreak{}Login}, directamente o por el \mbox{FA-03} del \mbox{CU-01}.}
  \textbf{Precondiciones} & \cuCelda{\begin{cureglas}
      \item \textbf{\mbox{PRE-01}} El cliente del SJM está instalado y en ejecución.
      \item \textbf{\mbox{PRE-02}} El Servidor de partidas está en ejecución y acepta conexiones TCP.
      \item \textbf{\mbox{PRE-03}} El Servidor de partidas tiene conexión disponible con la Base de datos.
      \item \textbf{\mbox{PRE-04}} El Jugador no tiene una sesión abierta en este cliente.
    \end{cureglas}} \\ \hline
  \textbf{Flujo normal} & \cuCelda{\begin{cuenum}[start=1]
      \item El SJM despliega la \texttt{GUI\_\allowbreak{}ForgotPassword}, solicitando el \texttt{correo} y mostrando las opciones ``Enviar'' y ``Cancelar''. \textit{(Ver \mbox{FA-08})}
      \item El Jugador captura su \texttt{correo} y da click en ``Enviar''.
      \item El SJM valida en el cliente que el campo esté completo y tenga formato de correo. \textit{(Ver \mbox{FA-01})}
      \item El SJM abre la conexión TCP con el Servidor de partidas y negocia la versión del protocolo. \textit{(Ver \mbox{EX-01}, \mbox{EX-02})}
      \item El SJM envía el mensaje \texttt{PASSWORD\_\allowbreak{}RESET\_\allowbreak{}REQUEST} con el correo capturado. \textit{(Ver \mbox{EX-06}, \mbox{EX-07})}
      \item El Servidor de partidas normaliza el correo a minúsculas y recupera de la Base de datos el \textit{Usuario} cuyo \texttt{correo} coincida. \textit{(Ver \mbox{EX-03})}
      \item El Servidor de partidas verifica que el \textit{Usuario} exista y que su \texttt{estado\_\allowbreak{}cuenta} no sea \texttt{ELIMINADA}. \textit{(Ver \mbox{FA-02})}
      \item El Servidor de partidas verifica que hayan transcurrido al menos cinco minutos desde \texttt{fecha\_\allowbreak{}ultimo\_\allowbreak{}envio\_\allowbreak{}recuperacion} (\mbox{RN-06}). \textit{(Ver \mbox{FA-03})}
    \end{cuenum}} \\
   & \cuCelda{\begin{cuenum}[start=9]
      \item El Servidor de partidas invalida los \textit{TokenRecuperacion} pendientes del \textit{Usuario} (\mbox{RN-05}).
      \item El Servidor de partidas genera un \textit{TokenRecuperacion} nuevo, guarda únicamente su resumen criptográfico y su fecha de expiración (\mbox{RN-04}, \mbox{RN-08}). \textit{(Ver \mbox{EX-03})}
      \item El Servidor de partidas actualiza \texttt{fecha\_\allowbreak{}ultimo\_\allowbreak{}envio\_\allowbreak{}recuperacion} del \textit{Usuario}.
      \item El Servidor de partidas solicita al Servicio de correo el envío del código en el \texttt{idioma\_\allowbreak{}preferido} del \textit{Usuario}. \textit{(Ver \mbox{EX-04})}
      \item El Servidor de partidas registra el intento en la \textit{BitacoraAcceso} con el resultado \texttt{RECUPERACION\_\allowbreak{}SOLICITADA}.
      \item El Servidor de partidas responde con \texttt{PASSWORD\_\allowbreak{}RESET\_\allowbreak{}SENT}.
      \item El SJM despliega la \texttt{GUI\_\allowbreak{}ResetPassword}, solicitando el código recibido, la contraseña nueva y su confirmación.
      \item El Jugador captura el código y la contraseña nueva con su confirmación, y da click en ``Cambiar contraseña''.
    \end{cuenum}} \\
   & \cuCelda{\begin{cuenum}[start=17]
      \item El SJM valida que la contraseña cumpla la política de \mbox{RN-03} y que coincida con su confirmación. \textit{(Ver \mbox{FA-04})}
      \item El SJM envía el mensaje \texttt{PASSWORD\_\allowbreak{}RESET\_\allowbreak{}CONFIRM} con el código y la contraseña nueva.
      \item El Servidor de partidas verifica que el \textit{TokenRecuperacion} exista, no esté usado y no haya expirado. \textit{(Ver \mbox{FA-05}, \mbox{FA-06})}
      \item El Servidor de partidas deriva el resumen de la contraseña nueva y comprueba que no coincida con \texttt{contrasena\_\allowbreak{}hash}. \textit{(Ver \mbox{FA-07})}
      \item El Servidor de partidas inicia una transacción sobre la Base de datos.
      \item El Servidor de partidas genera una \texttt{contrasena\_\allowbreak{}sal} nueva y deriva la \texttt{contrasena\_\allowbreak{}hash} de la contraseña recibida (\mbox{RN-02}). \textit{(Ver \mbox{EX-03})}
      \item El Servidor de partidas marca el \textit{TokenRecuperacion} como usado (\mbox{RN-04}).
      \item El Servidor de partidas restablece \texttt{intentos\_\allowbreak{}fallidos} a cero y deja \texttt{bloqueada\_\allowbreak{}hasta} en nulo.
    \end{cuenum}} \\
   & \cuCelda{\begin{cuenum}[start=25]
      \item El Servidor de partidas cierra todas las \textit{Sesion} abiertas del \textit{Usuario}, con \texttt{motivo\_\allowbreak{}cierre} igual a \texttt{CAMBIO\_\allowbreak{}CREDENCIALES} (\mbox{RN-07}).
      \item El Servidor de partidas confirma la transacción. \textit{(Ver \mbox{EX-05})}
      \item El Servidor de partidas registra el cambio en la \textit{BitacoraAcceso} con el resultado \texttt{RECUPERACION\_\allowbreak{}EXITOSA}.
      \item El Servidor de partidas responde con \texttt{PASSWORD\_\allowbreak{}RESET\_\allowbreak{}OK}.
      \item El SJM muestra la \texttt{GUI\_\allowbreak{}MessageSuccess} y vuelve a la \texttt{GUI\_\allowbreak{}Login}.
      \item Fin del caso de uso.
    \end{cuenum}} \\ \hline
  \textbf{Flujos alternos} & \cuCelda{\textbf{\mbox{FA-01}: El correo está incompleto o mal formado}\par
    \begin{cuenum}[start=1]
      \item El SJM muestra la \texttt{GUI\_\allowbreak{}MessageWarning} indicando qué le falta al campo y lo resalta.
      \item El flujo continúa en el paso 2 del FN. La validación es local y no llega al Servidor de partidas.
    \end{cuenum}} \\
   & \cuCelda{\textbf{\mbox{FA-02}: El correo no corresponde a ninguna cuenta}\par
    \begin{cuenum}[start=1]
      \item El Servidor de partidas no genera ningún token ni solicita ningún envío.
      \item El Servidor de partidas registra el intento en la \textit{BitacoraAcceso} con el resultado \texttt{RECUPERACION\_\allowbreak{}CORREO\_\allowbreak{}INEXISTENTE}, dejando \texttt{id\_\allowbreak{}usuario} en nulo.
      \item El Servidor de partidas espera el mismo tiempo que habría consumido el envío y responde con \texttt{PASSWORD\_\allowbreak{}RESET\_\allowbreak{}SENT}, idéntico al del FN (\mbox{RN-01}).
      \item El flujo continúa en el paso 15 del FN. El Jugador verá la pantalla del código y no podrá deducir si el correo existe.
    \end{cuenum}} \\
   & \cuCelda{\textbf{\mbox{FA-03}: Se solicita el envío antes del tiempo permitido}\par
    \begin{cuenum}[start=1]
      \item El Servidor de partidas responde con \texttt{PASSWORD\_\allowbreak{}RESET\_\allowbreak{}THROTTLED} y los segundos que faltan.
      \item El SJM los muestra en la \texttt{GUI\_\allowbreak{}MessageWarning}.
      \item El flujo continúa en el paso 2 del FN.
    \end{cuenum}} \\
   & \cuCelda{\textbf{\mbox{FA-04}: La contraseña nueva no cumple la política o no coincide}\par
    \begin{cuenum}[start=1]
      \item El SJM muestra la \texttt{GUI\_\allowbreak{}MessageWarning} explicando en una frase qué le falta conforme a \mbox{RN-03}, o bien que las dos capturas no coinciden, y limpia ambos campos.
      \item El flujo continúa en el paso 16 del FN.
    \end{cuenum}} \\
   & \cuCelda{\textbf{\mbox{FA-05}: El código es incorrecto}\par
    \begin{cuenum}[start=1]
      \item El Servidor de partidas incrementa en uno el atributo \texttt{intentos} del \textit{TokenRecuperacion}.
      \item El Servidor de partidas evalúa \texttt{intentos} contra \mbox{RN-09}. Si alcanza el máximo, invalida el token y lo registra en la \textit{BitacoraAcceso}.
      \item El Servidor de partidas responde con \texttt{PASSWORD\_\allowbreak{}RESET\_\allowbreak{}INVALID\_\allowbreak{}CODE} y los intentos restantes.
      \item Si se agotaron los intentos, el flujo continúa en el paso 1 del FN. En caso contrario, continúa en el paso 16.
    \end{cuenum}} \\
   & \cuCelda{\textbf{\mbox{FA-06}: El código expiró}\par
    \begin{cuenum}[start=1]
      \item El Servidor de partidas marca el \textit{TokenRecuperacion} como invalidado.
      \item El Servidor de partidas responde con \texttt{PASSWORD\_\allowbreak{}RESET\_\allowbreak{}EXPIRED}, que el SJM informa mediante la \texttt{GUI\_\allowbreak{}MessageWarning}.
      \item El flujo continúa en el paso 1 del FN, para que el Jugador solicite uno nuevo.
    \end{cuenum}} \\
   & \cuCelda{\textbf{\mbox{FA-07}: La contraseña nueva es igual a la anterior}\par
    \begin{cuenum}[start=1]
      \item El Servidor de partidas responde con \texttt{PASSWORD\_\allowbreak{}RESET\_\allowbreak{}SAME\_\allowbreak{}PASSWORD} sin consumir el token, para que el Jugador pueda reintentar con otra.
      \item El SJM muestra la \texttt{GUI\_\allowbreak{}MessageWarning} indicando que debe elegir una contraseña distinta.
      \item El flujo continúa en el paso 16 del FN.
    \end{cuenum}} \\
   & \cuCelda{\textbf{\mbox{FA-08}: Cancelar la recuperación}\par
    \begin{cuenum}[start=1]
      \item El Jugador selecciona ``Cancelar''.
      \item El SJM descarta lo capturado y despliega la \texttt{GUI\_\allowbreak{}Login}. Si ya se había generado un token, este sigue vigente hasta expirar.
      \item Fin del caso de uso.
    \end{cuenum}} \\ \hline
  \textbf{Excepciones} & \cuCelda{\textbf{\mbox{EX-01}: No se puede establecer la conexión con el servidor}\par
    \begin{cuenum}[start=1]
      \item El SJM detecta que la conexión TCP no pudo establecerse dentro del tiempo límite y registra el fallo en la bitácora local del cliente.
      \item El SJM muestra la \texttt{GUI\_\allowbreak{}MessageError} con las opciones ``Reintentar'' y ``Aceptar''.
      \item Si el Jugador selecciona ``Reintentar'', el flujo continúa en el paso 4 del FN. Si selecciona ``Aceptar'', continúa en el paso 2.
    \end{cuenum}} \\
   & \cuCelda{\textbf{\mbox{EX-02}: La versión del protocolo es incompatible}\par
    \begin{cuenum}[start=1]
      \item El Servidor de partidas responde con \texttt{PROTOCOL\_\allowbreak{}VERSION\_\allowbreak{}MISMATCH} y cierra la conexión.
      \item El SJM muestra la \texttt{GUI\_\allowbreak{}MessageError} indicando que debe actualizarse.
      \item Fin del caso de uso.
    \end{cuenum}} \\
   & \cuCelda{\textbf{\mbox{EX-03}: Fallo al consultar o actualizar la base de datos}\par
    \begin{cuenum}[start=1]
      \item El Servidor de partidas revierte la transacción en curso, si la hubiera, de modo que no quede una contraseña cambiada sin las sesiones cerradas.
      \item El Servidor de partidas registra la excepción en la bitácora del servidor con un identificador de incidencia, sin incluir la contraseña ni el código recibidos (\mbox{RN-02}).
      \item El Servidor de partidas responde con \texttt{SERVER\_\allowbreak{}ERROR} y el identificador de incidencia, y cierra la conexión.
      \item El SJM muestra la \texttt{GUI\_\allowbreak{}MessageError} con el identificador.
      \item Fin del caso de uso.
    \end{cuenum}} \\
   & \cuCelda{\textbf{\mbox{EX-04}: El servicio de correo no está disponible}\par
    \begin{cuenum}[start=1]
      \item El Servidor de partidas agota los reintentos previstos hacia el Servicio de correo sin obtener confirmación de entrega.
      \item El Servidor de partidas invalida el \textit{TokenRecuperacion} que había generado, para que no quede uno vigente que nadie recibió.
      \item El Servidor de partidas deja \texttt{fecha\_\allowbreak{}ultimo\_\allowbreak{}envio\_\allowbreak{}recuperacion} en su valor anterior, de modo que el Jugador pueda reintentar de inmediato sin esperar los cinco minutos de \mbox{RN-06}.
      \item El Servidor de partidas registra el fallo en la bitácora del servidor y responde con \texttt{EXTERNAL\_\allowbreak{}SERVICE\_\allowbreak{}UNAVAILABLE}.
      \item El SJM muestra la \texttt{GUI\_\allowbreak{}MessageError} indicando que no fue posible enviar el correo.
      \item Fin del caso de uso.
    \end{cuenum}} \\
   & \cuCelda{\textbf{\mbox{EX-05}: Fallo al confirmar la transacción}\par
    \begin{cuenum}[start=1]
      \item El Servidor de partidas no obtiene confirmación de que la transacción se haya escrito y la trata como fallida.
      \item El Servidor de partidas registra la excepción señalando que el estado del cambio es indeterminado, y responde con \texttt{SERVER\_\allowbreak{}ERROR}.
      \item El SJM muestra la \texttt{GUI\_\allowbreak{}MessageError} indicando que no fue posible confirmar el cambio y que intente iniciar sesión con la contraseña nueva antes de repetir el proceso.
      \item Fin del caso de uso.
    \end{cuenum}} \\
   & \cuCelda{\textbf{\mbox{EX-06}: Se agota el tiempo de espera de la respuesta}\par
    \begin{cuenum}[start=1]
      \item El SJM detecta que transcurrió el tiempo límite sin recibir un mensaje completo.
      \item El SJM cierra la conexión TCP y descarta el intercambio en curso, sin suponer que el servidor no lo procesó.
      \item El SJM registra el fallo en la bitácora local y lo informa mediante la \texttt{GUI\_\allowbreak{}MessageError}.
      \item El flujo continúa en el paso 2 del FN.
    \end{cuenum}} \\
   & \cuCelda{\textbf{\mbox{EX-07}: El mensaje recibido está malformado}\par
    \begin{cuenum}[start=1]
      \item El SJM no logra delimitar o deserializar el mensaje conforme al protocolo acordado.
      \item El SJM descarta el mensaje, cierra la conexión y registra en la bitácora local el contenido truncado.
      \item El SJM muestra la \texttt{GUI\_\allowbreak{}MessageError} indicando que la respuesta no pudo interpretarse.
      \item Fin del caso de uso.
    \end{cuenum}} \\ \hline
  \textbf{Postcondiciones} & \cuCelda{\begin{cureglas}
      \item \textbf{\mbox{POST-01}} Si el caso termina por el FN, el \textit{Usuario} tiene una \texttt{contrasena\_\allowbreak{}hash} y una \texttt{contrasena\_\allowbreak{}sal} nuevas, y la anterior ya no permite el acceso.
      \item \textbf{\mbox{POST-02}} Si el caso termina por el FN, el \textit{TokenRecuperacion} consumido queda marcado como usado y no puede volver a emplearse.
      \item \textbf{\mbox{POST-03}} Si el caso termina por el FN, ninguna \textit{Sesion} del \textit{Usuario} sigue abierta: todas tienen \texttt{fecha\_\allowbreak{}fin} y \texttt{motivo\_\allowbreak{}cierre} igual a \texttt{CAMBIO\_\allowbreak{}CREDENCIALES}.
      \item \textbf{\mbox{POST-04}} Si el caso termina por el FN, \texttt{intentos\_\allowbreak{}fallidos} vale cero y \texttt{bloqueada\_\allowbreak{}hasta} es nulo, de modo que un bloqueo temporal previo deja de impedir el acceso.
      \item \textbf{\mbox{POST-05}} Si el caso termina por el \mbox{FA-02}, no existe ningún \textit{TokenRecuperacion} nuevo y no se envió ningún correo, pero el Jugador recibió la misma respuesta que en el FN.
      \item \textbf{\mbox{POST-06}} Si el caso termina por el \mbox{EX-04}, existe la solicitud pero no hay token vigente ni correo enviado, y el Jugador puede reintentar sin esperar.
      \item \textbf{\mbox{POST-07}} En todos los casos en que el correo correspondía a una cuenta existente, el intento quedó registrado en la \textit{BitacoraAcceso} con su resultado, su fecha y hora y la dirección de red de origen.
    \end{cureglas}} \\ \hline
  \textbf{Reglas de negocio} & \cuCelda{\begin{cureglas}
      \item \textbf{\mbox{RN-01}} La respuesta del servidor es idéntica exista o no una cuenta con ese correo, y consume un tiempo equivalente. El proceso de recuperación no puede usarse para averiguar si una dirección está registrada.
      \item \textbf{\mbox{RN-02}} La contraseña nueva y el código no se almacenan, transmiten ni registran en claro. De ambos se guarda únicamente su resumen criptográfico con sal.
      \item \textbf{\mbox{RN-03}} La contraseña nueva cumple la misma política que el alta: al menos ocho caracteres y al menos tres de los cuatro grupos de minúsculas, mayúsculas, dígitos y signos. No puede contener el \texttt{nickname} ni el \texttt{correo}.
      \item \textbf{\mbox{RN-04}} El \textit{TokenRecuperacion} tiene una vigencia de treinta minutos y es de un solo uso.
      \item \textbf{\mbox{RN-05}} Solo puede existir un \textit{TokenRecuperacion} vigente por cuenta. Generar uno nuevo invalida el anterior.
      \item \textbf{\mbox{RN-06}} El correo de recuperación puede solicitarse a lo sumo una vez cada cinco minutos por cuenta.
      \item \textbf{\mbox{RN-07}} Cambiar la contraseña cierra todas las sesiones abiertas de la cuenta, en todos los dispositivos. Es la contrapartida de permitir varias sesiones simultáneas (\mbox{D-01}): si alguien entró a la cuenta, recuperar la contraseña debe echarlo.
    \end{cureglas}} \\
   & \cuCelda{\begin{cureglas}
      \item \textbf{\mbox{RN-08}} Del token se almacena su resumen, nunca su valor, por el mismo motivo que la contraseña.
      \item \textbf{\mbox{RN-09}} El código admite un máximo de tres intentos de verificación. Agotarlos lo invalida y obliga a solicitar uno nuevo.
      \item \textbf{\mbox{RN-10}} Una cuenta en \texttt{ELIMINADA} tiene el \texttt{correo} anonimizado, por lo que nunca coincide con el recibido y se trata como inexistente para efectos de \mbox{RN-01}.
    \end{cureglas}} \\ \hline
  \textbf{Observación} & \cuCelda{\textit{Sobre una tensión con el \mbox{CU-02}.}\par
    El \mbox{RN-01} oculta si un correo está registrado, pero el \mbox{FA-10} del \mbox{CU-02} lo dice abiertamente cuando alguien intenta darse de alta con un correo ya usado, y esa contradicción está aceptada a propósito en las observaciones de aquel caso. Se mantiene aquí la respuesta indistinguible de todos modos: cuesta lo mismo implementarla, y si más adelante se decide cerrar la fuga del registro, este caso ya no habría que tocarlo.} \\ \hline
  \textbf{Datos que toca} & \cuCelda{\begin{cudatos}
      \item \textbf{\textit{Usuario}:} lee \texttt{correo}, \texttt{estado\_\allowbreak{}cuenta}, \texttt{contrasena\_\allowbreak{}hash}, \texttt{idioma\_\allowbreak{}preferido} \textit{(pasos 6, 7, 20)}
      \item \textbf{\textit{Usuario}:} escribe \texttt{contrasena\_\allowbreak{}hash}, \texttt{contrasena\_\allowbreak{}sal}, \texttt{intentos\_\allowbreak{}fallidos}, \texttt{bloqueada\_\allowbreak{}hasta}, \texttt{fecha\_\allowbreak{}ultimo\_\allowbreak{}envio\_\allowbreak{}recuperacion} \textit{(pasos 11, 22, 24)}
      \item \textbf{\textit{TokenRecuperacion}:} crea, invalida y marca como usado \textit{(pasos 9, 10, 23)}
      \item \textbf{\textit{Sesion}:} cierra todas las abiertas \textit{(paso 25)}
      \item \textbf{\textit{BitacoraAcceso}:} inserta \textit{(pasos 13, 27)}
    \end{cudatos}} \\ \hline
\end{fichacu}
\clearpage
